% sdl_bgi-quickref.tex
%
% Guido Gonzato, PhD
% April 16, 2019

\documentclass[a4paper,11pt]{article}
\usepackage{graphicx}
\thispagestyle{empty}
\usepackage{color}
\usepackage{alltt}
\usepackage{framed}
\usepackage{setspace}
\usepackage[margin=3cm]{geometry}
\usepackage[colorlinks,urlcolor=blue]{hyperref}

\frenchspacing
\setlength{\parindent}{0pt}
\setlength{\parskip}{3pt}

\newcommand{\SDLbgi}{\texttt{SDL\_bgi}}

\newcommand{\version}{2.2.5}        % !!! <<<=== Change here !!!

\newcommand{\V}{\texttt{void}}      % void
\newcommand{\I}{\texttt{int}}       % int
\newcommand{\Ue}{\texttt{Uint8}}    % Uint8
\newcommand{\Ut}{\texttt{Uint32}}   % Uint32
\newcommand{\C}{\texttt{char *}}    % char *
\newcommand{\func}[1]{\textbf{\texttt{#1}}}  % function
\newcommand{\A}[1]{\emph{#1}}       % parameter
\newcommand{\F}[1]{\texttt{#1}}     % file
\newcommand{\T}[1]{\texttt{#1}}     % constant

\newcommand{\file}[1]{\texttt{#1}}

\definecolor{defaultshadecolor}{rgb}{0.85,0.85,1}
\definecolor{shadecolor}{rgb}{0.85,0.85,1}
\definecolor{gray}{rgb}{0.7,0.7,0.7}
\definecolor{charcolor}{rgb}{1,0.85,0.85}
\definecolor{lightred}{rgb}{1,0.7,0.7}
\definecolor{darkblue}{rgb}{0,0,0.5}
\definecolor{darkred}{rgb}{0.75,0,0}

\newenvironment{bgi}
{ % beg def
  \begin{snugshade}
}
{ % end def
  \end{snugshade}
}

% =====

\begin{document}

\title{Turtle Graphics Quick Reference}

\author{Guido Gonzato, PhD}

\date{\today}

\maketitle

\begin{center}
  \includegraphics[width=0.5\textwidth]{SDL_bgi_logo.png}  
\end{center}

% =====

\section{Turtle Graphics}

The \SDLbgi{} sources archive includes a \file{test/} directory that
contains test programs and also a turtle graphics implementation
(files \file{turtle.h} and \file{turtle.c}).

Turtle graphics is based on polar coordinates: lines are drawn
specifying the distance in pixels and an angle in degrees, called
\emph{heading} (0-359). Heading is the same as the bearing of a
compass: it starts from 0 (``North'') and increases clockwise. For
more information: \url{https://en.wikipedia.org/wiki/Turtle_graphics}

Some functions have aliases with the same name as their Python
counterparts; \url{https://docs.python.org/3/library/turtle.html}

% -----

\begin{bgi}
\V{} \func{back} (\I{} \A{px})
\end{bgi}

Moves the turtle backwards by \A{px} pixels.

Aliases: \func{backward()}, \func{bk()}

% -----

\begin{bgi}
\V{} \func{forwd} (\I{} \A{px})
\end{bgi}

Moves the turtle backwards by \A{px} pixels.

Aliases: \func{forward()}, \func{fd()}

% -----

\begin{bgi}
\V{} \func{turnleft} (\I{} \A{deg})
\end{bgi}

Turns the turtle left by \A{deg} degrees.

Aliases: \func{left()}, \func{lt()}

% -----

\begin{bgi}
\V{} \func{turnright} (\I{} \A{deg})
\end{bgi}

Turns the turtle right by \A{deg} degrees.

Aliases: \func{right()}, \func{rt()}

% -----

\begin{bgi}
\V{} \func{setposition} (\I{} \A{x} \I{} \A{y})
\end{bgi}

Moves the turtle to new coordinates \A{x}, \A{y}.

Aliases: \func{go\_to()}, \func{setpos()}

% -----

\begin{bgi}
\V{} \func{setx} (\I{} \A{x})
\end{bgi}

Sets the turtle's \A{x} coordinate.

% -----

\begin{bgi}
\V{} \func{sety} (\I{} \A{y})
\end{bgi}

Sets the turtle's \A{y} coordinate.

% -----

\begin{bgi}
\V{} \func{setheading} (\I{} \A{deg})
\end{bgi}

Sets the turtle heading to \A{deg} degrees.

Alias: \func{seth()}

% -----

\begin{bgi}
\V{} \func{home} (\V{})
\end{bgi}

Moves the turtle to the screen centre and sets the heading to 0.

% -----

\begin{bgi}
\I{} \func{* position} (\V{})
\end{bgi}

Returns a pointer to a 2-element integer array containing the turtle
coordinates.

% -----

\begin{bgi}
\I{} \func{xcor} (\V{})
\end{bgi}

Returns the turtle's \A{x} coordinate.

% -----

\begin{bgi}
\I{} \func{ycor} (\V{})
\end{bgi}

Returns the turtle's \A{y} coordinate.

% -----

\begin{bgi}
\I{} \func{heading} (\V{})
\end{bgi}

Returns the turtle's heading in degrees.

% -----

\begin{bgi}
\V{} \func{savestate} (\V{})
\end{bgi}

Saves the turtle state, i.e.\ its coordinates and heading.

TODO: make this a stack?

% -----

\begin{bgi}
\V{} \func{restorestate} (\V{})
\end{bgi}

Restore previously saved state. If no state has been saved, this
function is equivalent to \func{home()}.

% -----

\begin{bgi}
\V{} \func{pendown} (\V{})
\end{bgi}

Activates drawing (puts the pen down).

Aliases: \func{pd()}, \func{down()}

% -----

\begin{bgi}
\V{} \func{penup} (\V{})
\end{bgi}

Suspends drawing (lifts the pen down).

Aliases: \func{pu()}, \func{up()}

% -----

\begin{bgi}
\I{} \func{isdown} (\V{})
\end{bgi}

Returns 1 if the pen is down.

% -----

\begin{bgi}
\V{} \func{hideturtle} (\V{})
\end{bgi}

Hides the turtle.

% -----

\begin{bgi}
\V{} \func{showturtle} (\V{})
\end{bgi}

Draws the turtle.

% -----

\begin{bgi}
\I{} \func{isvisible} (\V{})
\end{bgi}

Returns 1 if the turtle is visible.

% -----

\begin{bgi}
\I{} \func{turtleshape} (\I{} \A{shape})
\end{bgi}

Sets the turtle shape; available shapes are \T{T\_CIRCLE} or
\T{T\_TRIANGLE}.

% -----

\begin{bgi}
\I{} \func{turtlesize} (\I{} \A{px})
\end{bgi}

Sets the turtle size (radius or side) to \A{px} pixels.

% -----

\begin{bgi}
\V{} \func{wrap} (\V{})
\end{bgi}

Wraps around the screen.

% -----

\begin{bgi}
\V{} \func{wrap} (\V{})
\end{bgi}

Does not wrap around the screen.

% =====

\bigskip

\rule{\linewidth}{0.2pt}

This document is a free manual, released under the GNU Free
Documentation License (FDL) v. 1.3 or later.

\end{document}

% ----- end of file sdl_bgi-quickref.tex
